% ============================================================
%  Math Notes Template
% ============================================================
\documentclass[11pt, a4paper]{article}

% --- Mathematics (must load before newtxmath) ---
\usepackage{amsmath, amssymb, amsthm, mathtools}

% --- Fonts & Encoding (XeLaTeX) ---
\usepackage{fontspec}
\usepackage{libertine}                    % Linux Libertine text + Biolinum sans
\usepackage[libertine]{newtxmath}         % Matching math font (loads after amsthm)
\usepackage{inconsolata}                  % Monospace font
\usepackage{microtype}                    % Improved typography

% --- Page Layout ---
\usepackage[
  margin=2.5cm,
  headheight=14pt
]{geometry}

% --- Line Spacing ---
\usepackage{setspace}
\onehalfspacing

% --- Paragraph Spacing ---
\usepackage{parskip}

% --- Colours ---
\usepackage[dvipsnames]{xcolor}
\definecolor{defblue}{HTML}{1a5276}
\definecolor{thmgreen}{HTML}{196f3d}
\definecolor{exred}{HTML}{922b21}
\definecolor{linkblue}{HTML}{2471a3}

% --- TikZ ---
\usepackage{tikz}
\usetikzlibrary{
  arrows.meta,
  calc,
  decorations.markings,
  patterns,
  positioning,
  shapes.geometric,
  cd                                      % commutative diagrams
}

% --- Theorem Environments ---
\usepackage{thmtools}

\declaretheoremstyle[
  headfont=\bfseries\color{thmgreen},
  bodyfont=\normalfont,
  spaceabove=10pt,
  spacebelow=10pt,
]{thmstyle}

\declaretheoremstyle[
  headfont=\bfseries\color{defblue},
  bodyfont=\normalfont,
  spaceabove=10pt,
  spacebelow=10pt,
]{defstyle}

\declaretheoremstyle[
  headfont=\bfseries\color{exred},
  bodyfont=\normalfont,
  spaceabove=10pt,
  spacebelow=10pt,
]{exstyle}

\declaretheorem[style=thmstyle, numberwithin=section, name=Theorem]{theorem}
\declaretheorem[style=thmstyle, sibling=theorem, name=Lemma]{lemma}
\declaretheorem[style=thmstyle, sibling=theorem, name=Proposition]{proposition}
\declaretheorem[style=thmstyle, sibling=theorem, name=Corollary]{corollary}
\declaretheorem[style=defstyle, sibling=theorem, name=Definition]{definition}
\declaretheorem[style=exstyle,  sibling=theorem, name=Example]{example}
\declaretheorem[style=exstyle,  numbered=no,      name=Remark]{remark}

% --- Hyperlinks & Cross-references ---
\usepackage[
  colorlinks=true,
  linkcolor=linkblue,
  citecolor=linkblue,
  urlcolor=linkblue
]{hyperref}
\usepackage{cleveref}

% --- Headers & Footers ---
\usepackage{fancyhdr}
\pagestyle{fancy}
\fancyhf{}
\fancyhead[L]{\textsc{\nouppercase{\leftmark}}}
\fancyhead[R]{\thepage}
\renewcommand{\headrulewidth}{0.4pt}

% --- Enumeration ---
\usepackage{enumitem}

% --- Bibliography (optional, uncomment if needed) ---
% \usepackage[style=alphabetic, backend=biber]{biblatex}
% \addbibresource{refs.bib}

% ============================================================
%  Fill these in for each set of notes
% ============================================================
\newcommand{\notestitle}{Algebraic Geometry I: Schemes}
\newcommand{\notesauthor}{B. Arm}
\newcommand{\notesdate}{\today}

\title{\notestitle}
\author{\notesauthor}
\date{\notesdate}

% ============================================================
\begin{document}

\maketitle
\tableofcontents
\newpage

% ============================================================
%  Chapter 1: Prevarieties
% ============================================================
\section{Prevarieties}

\subsection{Affine Algebraic Sets}

\subsection{Affine Algebraic Sets as Spaces with Functions}

\subsection{Prevarieties}

\subsection{Projective Varieties}

% ============================================================
%  Chapter 2: Spectrum of a Ring
% ============================================================
\section{Spectrum of a Ring}

\subsection{Spectrum of a Ring as a Topological Space}

\subsection{Excursion: Sheaves}

\subsection{Spectrum of a Ring as a Locally Ringed Space}

% ============================================================
%  Chapter 3: Schemes
% ============================================================
\section{Schemes}

\subsection{Schemes}

\subsection{Examples of Schemes}

\subsection{Basic Properties of Schemes and Morphisms of Schemes}

\subsection{Prevarieties as Schemes}

\subsection{Subschemes and Immersions}

% ============================================================
%  Chapter 4: Fiber Products
% ============================================================
\section{Fiber Products}

\subsection{Schemes as Functors}

\subsection{Fiber Products of Schemes}

\subsection{Base Change, Fibers of a Morphism}

% ============================================================
%  Chapter 5: Schemes over Fields
% ============================================================
\section{Schemes over Fields}

\subsection{Schemes over a Field Which Is Not Algebraically Closed}

\subsection{Dimension of Schemes over a Field}

\subsection{Schemes over Fields and Extensions of the Base Field}

\subsection{Intersections of Plane Curves}

% ============================================================
%  Chapter 6: Local Properties of Schemes
% ============================================================
\section{Local Properties of Schemes}

\subsection{The Tangent Space}

\subsection{Smooth Morphisms}

\subsection{Regular Schemes}

\subsection{Normal Schemes}

% ============================================================
%  Chapter 7: Quasi-coherent Modules
% ============================================================
\section{Quasi-coherent Modules}

\subsection{Excursion: $\mathcal{O}_X$-Modules}

\subsection{Quasi-coherent Modules on a Scheme}

\subsection{Properties of Quasi-coherent Modules}

% ============================================================
%  Chapter 8: Representable Functors
% ============================================================
\section{Representable Functors}

\subsection{Representable Functors}

\subsection{The Example of the Grassmannian}

\subsection{Brauer--Severi Schemes}

% ============================================================
%  Chapter 9: Separated Morphisms
% ============================================================
\section{Separated Morphisms}

\subsection{Diagonal of Scheme Morphisms and Separated Morphisms}

\subsection{Rational Maps and Function Fields}

% ============================================================
%  Chapter 10: Finiteness Conditions
% ============================================================
\section{Finiteness Conditions}

\subsection{Finiteness Conditions (Noetherian Case)}

\subsection{Finiteness Conditions in the Non-Noetherian Case}

\subsection{Schemes over Inductive Limits of Rings}

\subsection{Constructible Properties}

% ============================================================
%  Chapter 11: Vector Bundles
% ============================================================
\section{Vector Bundles}

\subsection{Vector Bundles and Locally Free Modules}

\subsection{Flattening Stratification for Modules}

\subsection{Divisors}

\subsection{Vector Bundles on $\mathbb{P}^1$}

% ============================================================
%  Chapter 12: Affine and Proper Morphisms
% ============================================================
\section{Affine and Proper Morphisms}

\subsection{Affine Morphisms}

\subsection{Finite and Quasi-finite Morphisms}

\subsection{Serre's and Chevalley's Criteria to Be Affine}

\subsection{Normalization}

\subsection{Proper Morphisms}

\subsection{Zariski's Main Theorem}

% ============================================================
%  Chapter 13: Projective Morphisms
% ============================================================
\section{Projective Morphisms}

\subsection{Projective Spectrum of a Graded Algebra}

\subsection{Embeddings into Projective Space}

\subsection{Blowing-up}

% ============================================================
%  Chapter 14: Flat Morphisms and Dimension
% ============================================================
\section{Flat Morphisms and Dimension}

\subsection{Flat Morphisms}

\subsection{Properties of Flat Morphisms}

\subsection{Faithfully Flat Descent}

\subsection{Dimension and Fibers of Morphisms}

\subsection{Dimension and Regularity Conditions}

\subsection{Hilbert Schemes}

% ============================================================
%  Chapter 15: One-dimensional Schemes
% ============================================================
\section{One-dimensional Schemes}

\subsection{Morphisms into and from One-dimensional Schemes}

\subsection{Valuative Criteria}

\subsection{Curves over Fields}

\subsection{Divisors on Curves}

% ============================================================
%  Chapter 16: Examples
% ============================================================
\section{Examples}

\subsection{Determinantal Varieties}

\subsection{Cubic Surfaces and a Hilbert Modular Surface}

\subsection{Cyclic Quotient Singularities}

\subsection{Abelian Varieties}

% ============================================================
%  Appendices
% ============================================================
\section*{Appendix A: The Language of Categories}
\addcontentsline{toc}{section}{Appendix A: The Language of Categories}

\section*{Appendix B: Commutative Algebra}
\addcontentsline{toc}{section}{Appendix B: Commutative Algebra}

% \printbibliography

\end{document}
